\section*{الفصل \ref{ch:b3:solid-state} --- كيمياء الحالة الصلبة: الحزم والعيوب وأشباه الموصلات}

\begin{solution}{exo:b3:solid-state:1}
$(E - \alpha)/|\beta| = -2\cos(k\pi/7)$: $-1.802$، $-1.247$، $-0.445$، $+0.445$، $+1.247$، $+1.802$. وهي تمتد على
$3.604|\beta|$، أي 90\,\% من عرض الحزمة $4|\beta|$ للسلسلة اللانهائية منذ الآن.
\end{solution}

\begin{solution}{exo:b3:solid-state:2}
$2N$ (عدد $N$ من المدارات، يتسع كل منها لإلكترونين). الصوديوم، بإلكترون 3s واحد لكل ذرة، يملأ
نصف حزمته s: فهو فلز. وللمغنيزيوم إلكترونان وكان سيملأ الحزمة s تمامًا،
لكن الحزمتين 3s و3p متداخلتان، فتنتمي أعلى المستويات المشغولة إلى
حزمة مركّبة مملوءة جزئيًا: فهو فلز أيضًا.
\end{solution}

\begin{solution}{exo:b3:solid-state:3}
$\lambda \approx \qty{1239.8}{eV.nm}/E_g$: فوسفيد الغاليوم \qty{549}{nm} (أخضر)؛ ونتريد الغاليوم \qty{365}{nm}
(فوق البنفسجي القريب). وتستعمل الصمامات الزرقاء نتريد الغاليوم مخلوطًا بالإنديوم، الذي
فجوته أصغر.
\end{solution}

\begin{solution}{exo:b3:solid-state:4}
$\ce{CaCl2} \xrightarrow{\ce{NaCl}} \mathrm{Ca_{Na}^{\bullet} + V_{Na}' + 2\,Cl_{Cl}^{\times}}$: موقعان كاتيونيان وموقعان أنيونيان، كما في NaCl؛ وCa واحد 
وCl اثنان؛ والشحنة $+1 - 1 = 0$. $\ce{Y2O3} \xrightarrow{\ce{ZrO2}} \mathrm{2\,Y_{Zr}' + 3\,O_O^{\times} + V_O^{\bullet\bullet}}$: موقعان كاتيونيان وأربعة مواقع أكسجين؛ وY اثنان
وO ثلاثة؛ والشحنة $-2 + 2 = 0$.
\end{solution}

\begin{solution}{exo:b3:solid-state:5}
$T^{3/2} = 5196$ عند \qty{300}{K}، $kT = \qty{0.02585}{eV}$. السيليكون: $\sqrt{N_cN_v} = \qty{2.42e19}{cm^{-3}}$، $\eu^{-1.125/0.05170} = \num{3.55e-10}$،
$n_i = \qty{8.6e9}{cm^{-3}}$. الجرمانيوم: $\sqrt{N_cN_v} = \qty{7.16e18}{cm^{-3}}$، $\eu^{-0.661/0.05170} = \num{2.80e-6}$، $n_i = \qty{2.0e13}{cm^{-3}}$. والنسبة
$\num{4.3e-4}$: فالفجوة الأصغر للجرمانيوم تمنحه حاملات أكثر بأكثر من ألفي
مرة.
\end{solution}

\begin{solution}{exo:b3:solid-state:6}
$n = N_D = \qty{1.0e16}{cm^{-3}}$؛ $p = n_i^2/n = \num{1.0e20}/\num{1.0e16} = \qty{1.0e4}{cm^{-3}}$.
\end{solution}

\begin{solution}{exo:b3:solid-state:7}
$E_F - E_i = kT\ln(n/n_i) = \qty{0.02585}{eV} \times \ln(10^6) = \qty{0.36}{eV}$.
\end{solution}

\begin{solution}{exo:b3:solid-state:8}
$kT = \qty{0.06894}{eV}$ عند \qty{800}{K}؛ $n/N = \eu^{-2.3/0.1379} = \num{5.7e-8}$؛ وفي مول $\num{6.02e23} \times \num{5.7e-8} = \num{3.4e16}$
زوج.
\end{solution}

\begin{solution}{exo:b3:solid-state:9}
كتلة الخلية $\rho a^3 = \qty{5.70}{g.cm^{-3}} \times (\qty{4.300e-8}{cm})^3 = \qty{4.532e-22}{g}$، أي
$\qty{4.532e-22}{g} \times N_A = \qty{272.9}{g}$ لكل مول من الخلايا، أو \qty{68.23}{g} لكل أكسجين. ثم
تعطي $(1 - x) \times 55.8 + 16.0 = 68.23$ القيمة $1 - x = 0.936$: $x = 0.064$، $\mathrm{Fe_{0.936}O}$.
\end{solution}

\begin{solution}{exo:b3:solid-state:10}
$\ln(\sigma T)$: $-3.51$، $-1.30$، $0.36$ عند $1/T = 3.333$، $2.857$، $\qty{2.500e-3}{K^{-1}}$؛ والميل بين النقطتين الطرفيتين
$(0.36 + 3.51)/(-\qty{8.33e-4}{K^{-1}}) = -\qty{4.65e3}{K}$ (والنقطة الوسطى تقع على المستقيم)، ومنه
$E_a = \qty{4.65e3}{K} \times \qty{8.617e-5}{eV.K^{-1}} = \qty{0.40}{eV}$.
\end{solution}

\begin{solution}{exo:b3:solid-state:11}
$\tfrac12\ce{O2} \rightleftharpoons \mathrm{O_O^{\times} + V_M' + h^{\bullet}}$، $K = [\mathrm{V_M'}][\mathrm h^\bullet]/p(\ce{O2})^{1/2}$. وحصيلة الشحنة $[\mathrm h^\bullet] = [\mathrm{V_M'}]$، ومنه
$[\mathrm h^\bullet]^2 = K\,p(\ce{O2})^{1/2}$ و$[\mathrm h^\bullet] \propto p(\ce{O2})^{1/4}$؛ والناقلية تتبع \omterm{def:b3:solid-state:carriers}{الثقوب}.
\end{solution}

\begin{solution}{exo:b3:solid-state:12}
$n/N = \sqrt{N_i/N}\,\eu^{-\Delta H_F/2kT} = \sqrt 2\,\eu^{-1.4/(2 \times 0.05170)} = 1.414 \times \num{1.32e-6} = \num{1.9e-6}$.
\end{solution}

\begin{solution}{pb:b3:solid-state:1}
\textbf{1.} $\ce{Y2O3} \xrightarrow{\ce{ZrO2}} \mathrm{2\,Y_{Zr}' + 3\,O_O^{\times} + V_O^{\bullet\bullet}}$.

\textbf{2.} المواقع: موقعان كاتيونيان وأربعة مواقع أكسجين (ثلاثة ممتلئة، وواحد فارغ)،
أي النسبة 1\,:\,2 في $\ce{ZrO2}$. المادة: 2 Y و3 O في كل طرف. الشحنة: $2 \times (-1) + 2 = 0$.

\textbf{3.} فراغ واحد لكل أيوني إيتريوم: أي نصف فراغ لكل إيتريوم.

\textbf{4.} لكل $(\ce{ZrO2})_{0.92}(\ce{Y2O3})_{0.08}$ يوجد 0.92 Zr و0.16 Y، أي 1.08 كاتيون: $x = 0.16/1.08 = 0.148$،
$\mathrm{Zr_{0.852}Y_{0.148}O_{1.926}}$.

\textbf{5.} $(x/2)/2 = 0.0370$: أي إن 3.7\,\% من مواقع الأكسجين فارغة.

\textbf{6.} $8 \times 0.0370 = 0.296$ فراغ في الخلية.

\textbf{7.} $a^3 = (\qty{5.14e-8}{cm})^3 = \qty{1.358e-22}{cm^3}$: $0.296/\qty{1.358e-22}{cm^3} = \qty{2.2e21}{cm^{-3}}$.

\textbf{8.} يجب أن يقفز كل أيون أكسيد فوق حاجز قدره \qty{1.0}{eV} إلى فراغ
مجاور؛ ونسبة المحاولات الناجحة، $\eu^{-E_a/kT}$، تكبر بسرعة كبيرة مع $T$.

\textbf{9.} عند \qty{1073.15}{K}، $kT = \qty{0.09248}{eV}$: $A = \sigma T\eu^{E_a/kT} = 0.030 \times 1073.15 \times \eu^{10.81} = \qty{1.6e6}{S.K.cm^{-1}}$.

\textbf{10.} عند \qty{973.15}{K}: $\sigma = (A/T)\eu^{-11.92} = \qty{0.011}{S.cm^{-1}}$.

\textbf{11.} عند \qty{573.15}{K}: $\sigma = (A/T)\eu^{-20.25} = \qty{4.5e-6}{S.cm^{-1}}$.

\textbf{12.} $R = L/(\sigma S) = \qty{0.10}{cm}/(\sigma \times \qty{0.20}{cm^2})$: \qty{46}{\ohm} عند \qty{700}{\celsius}، و$\qty{1.1e5}{\ohm}$ عند \qty{300}{\celsius}.

\textbf{13.} يقتضي $R < \qty{1}{k\ohm}$ أن يكون $\sigma > \qty{5.0e-4}{S.cm^{-1}}$؛ وحل $(A/T)\eu^{-E_a/kT} = \num{5.0e-4}$ (بالتجريب)
يعطي $T \approx \qty{760}{K}$، أي نحو \qty{490}{\celsius}.

\textbf{14.} يكون العادم باردًا بعد التشغيل وعند الحمل المنخفض؛ ويوصل السخّان
الإلكتروليت إلى درجة حرارة عمله في ثوانٍ، بحيث يعمل التحكم
منذ البداية، حين يُطلق معظم الملوثات.

\textbf{15.} $\ce{O2 + 4e- -> 2O^2-}$؛ وفي \omterm{def:b3:solid-state:kroger-vink}{ترميز كروغر--فينك} $\ce{O2} + \mathrm{2\,V_O^{\bullet\bullet} + 4\,e'} \longrightarrow \mathrm{2\,O_O^{\times}}$.

\textbf{16.} يُختزل الأكسجين في جانب الهواء وتتأكسد أيونات الأكسيد عائدةً
إلى أكسجين في جانب العادم: فتنقل الخلية $\ce{O2}$ من الهواء إلى العادم، مع
$\Delta_rG = RT\ln(p_{\mathrm{exh}}/p_{\mathrm{air}})$ لكل مول من $\ce{O2}$ وأربعة إلكترونات: $E = -\Delta_rG/4F = (RT/4F)\ln(p_{\mathrm{air}}/p_{\mathrm{exh}})$.

\textbf{17.} $RT/4F = 8.3145 \times 973.15/(4 \times 96485) = \qty{0.02096}{V}$.

\textbf{18.} صفر: فالضغطان متساويان.

\textbf{19.} $(RT/4F)\ln 10 = \qty{48.3}{mV}$ لكل عُشر.

\textbf{20.} $\qty{0.02096}{V} \times \ln(0.2095/0.010) = \qty{0.02096}{V} \times 3.04 = \qty{0.064}{V}$.

\textbf{21.} $\qty{0.02096}{V} \times \ln(0.2095/\num{1.0e-20}) = \qty{0.02096}{V} \times 44.49 = \qty{0.93}{V}$.

\textbf{22.} يحوي العادم الغني أحادي أكسيد الكربون والهيدروجين والهيدروكربونات غير
المحترقة؛ وعلى قطب البلاتين يتفاعل معها القليل الباقي من الأكسجين،
ويتحدد $p(\ce{O2})$ بالتوازنين \ce{2CO + O2 <=> 2CO2} و\ce{2H2 + O2 <=> 2H2O}، اللذين يتركان
قيمة ضئيلة جدًا عند \qty{700}{\celsius}.

\textbf{23.} $p = 0.2095\,\eu^{-0.45/0.02096} = \qty{1.0e-10}{bar}$.

\textbf{24.} حول النسبة الستوكيومترية، ينتقل العادم من فائض طفيف
من الأكسجين إلى فائض طفيف من CO والهيدروجين: فينخفض $p(\ce{O2})$ بنحو ثماني عشرة
رتبة عشرية من أجل تغير صغير في الوقود، ويقفز الجهد من أقل من
\qty{0.1}{V} إلى نحو \qty{0.9}{V}. ويضيف التحكم في المحرك وقودًا حين يكون الجهد منخفضًا 
وينقصه حين يكون مرتفعًا، فيُبقي الخليط عند النقطة الستوكيومترية حيث يحوّل
المحوّل الحفزي الثلاثي CO والهيدروكربونات وأكاسيد النيتروجين في آن واحد.

\textbf{25.} في العادم الغني عند \qty{700}{\celsius} يعطي المستشعر نحو \qty{0.93}{V}.
\end{solution}
